\documentclass[11pt,a4paper]{article}
\usepackage[utf8]{inputenc}
\usepackage[T1]{fontenc}
\usepackage[margin=0.72in]{geometry}
\usepackage{enumitem}
\usepackage{titlesec}
\usepackage{xcolor}
\usepackage{hyperref}
\usepackage{microtype}

\definecolor{primarycolor}{RGB}{33,37,41}
\definecolor{secondarycolor}{RGB}{0,92,153}
\definecolor{accentcolor}{RGB}{90,99,108}
\hypersetup{colorlinks=true,urlcolor=secondarycolor,linkcolor=secondarycolor}
\pagestyle{empty}
\setlist[itemize]{leftmargin=1.4em,itemsep=0.2em,topsep=0.25em}
\titleformat{\section}{\large\bfseries\color{primarycolor}}{}{0em}{}
\titlespacing*{\section}{0pt}{1.1em}{0.45em}
\setlength{\parindent}{0pt}
\setlength{\parskip}{0.35em}

\begin{document}

\begin{center}
{\LARGE\bfseries\color{primarycolor} MARTA KOWAL-BRANDT}\\
QA Automation Engineer | CI/CD \& Platform Engineering\\
Wroclaw, Poland | marta.kowalbrandt@example.com | +48 555 303 303 | gitlab.com/mkb-example
\end{center}

\section*{PROFESSIONAL SUMMARY}
QA automation engineer with 7 years of experience, including hands-on ownership of CI/CD pipelines and test infrastructure for three product teams. Built GitLab CI workflows and Helm-based ephemeral environments on a shared Kubernetes cluster; personal projects include a k3s GitOps homelab and Terraform-provisioned Hetzner VPS infrastructure. Seeking to move into a DevOps or platform engineering role.

\section*{SKILLS}
Infrastructure and delivery: GitLab CI, Jenkins, Docker, Docker Compose, Helm, Kubernetes (namespace-level access at work; k3s in a personal project)\\
Infrastructure as code and cloud learning: Terraform (personal Hetzner VPS project), Ansible (basic), AWS (personal learning: EC2 and S3 basics)\\
Observability: Prometheus, Grafana, Loki; Grafana dashboards for test metrics\\
Scripting and data: Python, Bash, SQL, PostgreSQL\\
Systems and testing: Linux, Git, pytest, Playwright, Selenium, Postman, Java (test code)

\section*{EXPERIENCE}
\textbf{QA Automation Engineer (Test Infrastructure) | Odrava Health Software, Wroclaw} \hfill \textit{2021–Present}
\begin{itemize}
\item Own GitLab CI pipelines for three product teams (about 25 developers), covering build, unit, integration and end-to-end test stages, and artifact publishing.
\item Reduced average pipeline time from 38 to 14 minutes by parallelising test jobs, caching dependencies and separating end-to-end tests into their own stage; monthly CI runner cost fell by about 30\%.
\item Built merge-request ephemeral test environments using Docker Compose and later Helm charts on the company’s shared Kubernetes cluster; environments include the application, PostgreSQL and seeded data and are automatically destroyed after merge or 48 hours of inactivity.
\item Maintain a pytest and Playwright end-to-end suite of about 900 tests; reduced the flaky-test rate from 7\% to below 1\%.
\item Wrote Bash and Python tooling for PostgreSQL test-data seeding and an internal guide to the CI system used in onboarding.
\end{itemize}

\textbf{QA Engineer | Brightmill Software House, Wroclaw} \hfill \textit{2019–2021}
\begin{itemize}
\item Tested web and mobile applications for client projects using manual and automated testing.
\item Introduced a Selenium and Java automated regression suite for the largest client project and configured Jenkins jobs for nightly runs.
\end{itemize}

\textbf{Customer Support Specialist (2nd line) | Brightmill Software House} \hfill \textit{2018–2019}
\begin{itemize}
\item Investigated customer issues in logs and databases and reproduced bugs for developers.
\end{itemize}

\textbf{Homelab k3s Cluster with GitOps | Personal} \hfill \textit{2024–2025}
\begin{itemize}
\item Built a three-node Raspberry Pi 5 k3s cluster and deployed applications from a Git repository using Argo CD.
\item Added Prometheus, Grafana and Loki for monitoring, Traefik ingress and cert-manager for TLS; provisioned infrastructure for a small Hetzner VPS with Terraform.
\end{itemize}

\section*{SELECTED PROJECTS}
\textbf{Merge-Request Ephemeral Environments | Odrava Health} \hfill \textit{2023–2024}
\begin{itemize}
\item Implemented short-lived environments for each merge request, deployed via Helm to the shared Kubernetes cluster from GitLab CI, with automatic cleanup after merge or 48 hours of inactivity.
\item Enabled QA and product teams to review features before merge and run end-to-end tests against the environments.
\end{itemize}

\textbf{CI Pipeline Speed-up | Odrava Health} \hfill \textit{2022}
\begin{itemize}
\item Profiled pipelines for three teams, parallelised test jobs, added dependency caching and split end-to-end tests into a separate stage.
\item Reduced average pipeline time from 38 to 14 minutes and monthly CI runner cost by about 30\%.
\end{itemize}

\section*{EDUCATION}
BEng Computer Engineering | Silesian Lowlands University of Technology | 2018

\section*{CERTIFICATION AND DEVELOPMENT}
ISTQB Foundation Level | 2019\\
Certified Kubernetes Administrator (CKA) preparation; exam booked for January 2027 (not yet certified)

\section*{LANGUAGES}
Polish (native), English (C1), German (A2)

\end{document}